% Options for packages loaded elsewhere
\PassOptionsToPackage{unicode}{hyperref}
\PassOptionsToPackage{hyphens}{url}
\documentclass[
  11pt,
  a4paper,
]{article}
\usepackage{xcolor}
\usepackage[margin=18mm]{geometry}
\usepackage{amsmath,amssymb}
\setcounter{secnumdepth}{-\maxdimen} % remove section numbering
\usepackage{iftex}
\ifPDFTeX
  \usepackage[T1]{fontenc}
  \usepackage[utf8]{inputenc}
  \usepackage{textcomp} % provide euro and other symbols
\else % if luatex or xetex
  \usepackage{unicode-math} % this also loads fontspec
  \defaultfontfeatures{Scale=MatchLowercase}
  \defaultfontfeatures[\rmfamily]{Ligatures=TeX,Scale=1}
\fi
\ifPDFTeX\else
  % xetex/luatex font selection
    \setmainfont[]{Noto Serif}
    \setsansfont[]{Noto Sans}
    \setmonofont[]{Noto Sans Mono}
  \setmathfont[]{Noto Sans Math}
\fi
% Use upquote if available, for straight quotes in verbatim environments
\IfFileExists{upquote.sty}{\usepackage{upquote}}{}
\IfFileExists{microtype.sty}{% use microtype if available
  \usepackage[]{microtype}
  \UseMicrotypeSet[protrusion]{basicmath} % disable protrusion for tt fonts
}{}
\makeatletter
\@ifundefined{KOMAClassName}{% if non-KOMA class
  \IfFileExists{parskip.sty}{%
    \usepackage{parskip}
  }{% else
    \setlength{\parindent}{0pt}
    \setlength{\parskip}{6pt plus 2pt minus 1pt}}
}{% if KOMA class
  \KOMAoptions{parskip=half}}
\makeatother
\ifLuaTeX
\usepackage[bidi=basic]{babel}
\else
\usepackage[bidi=default]{babel}
\fi
\babelprovide[main,import]{russian}
\ifPDFTeX
\else
\babelfont{rm}[]{Noto Serif}
\fi
% get rid of language-specific shorthands (see #6817):
\let\LanguageShortHands\languageshorthands
\def\languageshorthands#1{}
\setlength{\emergencystretch}{3em} % prevent overfull lines
\providecommand{\tightlist}{%
  \setlength{\itemsep}{0pt}\setlength{\parskip}{0pt}}
\usepackage{bookmark}
\IfFileExists{xurl.sty}{\usepackage{xurl}}{} % add URL line breaks if available
\urlstyle{same}
\hypersetup{
  pdftitle={Билеты по дискретным случайным процессам},
  pdflang={ru-RU},
  hidelinks,
  pdfcreator={LaTeX via pandoc}}

\title{Билеты по дискретным случайным процессам}
\author{}
\date{}

\begin{document}
\maketitle

\section{\texorpdfstring{Билет 03. Продолжение меры с полуалгебры на
алгебру и \(\sigma\)-алгебру. Мера
Лебега}{Билет 03. Продолжение меры с полуалгебры на алгебру и \textbackslash sigma-алгебру. Мера Лебега}}\label{ux431ux438ux43bux435ux442-03.-ux43fux440ux43eux434ux43eux43bux436ux435ux43dux438ux435-ux43cux435ux440ux44b-ux441-ux43fux43eux43bux443ux430ux43bux433ux435ux431ux440ux44b-ux43dux430-ux430ux43bux433ux435ux431ux440ux443-ux438-sigma-ux430ux43bux433ux435ux431ux440ux443.-ux43cux435ux440ux430-ux43bux435ux431ux435ux433ux430}

Источники: Ширяев 1, гл. II, §3; лекции ДСП.

\subsection{1. Короткий
ответ}\label{ux43aux43eux440ux43eux442ux43aux438ux439-ux43eux442ux432ux435ux442}

Меру часто задают сначала на простых множествах, где ее значение
очевидно, а затем продолжают на более богатые классы:

\[
\mathcal S\longrightarrow\mathcal A(\mathcal S)\longrightarrow\sigma(\mathcal S).
\]

Здесь \(\mathcal S\) - полуалгебра, \(\mathcal A(\mathcal S)\) - алгебра
конечных дизъюнктных объединений элементов \(\mathcal S\), а
\(\sigma(\mathcal S)\) - порожденная \(\sigma\)-алгебра.

Если

\[
A=\bigsqcup_{k=1}^n S_k,\qquad S_k\in\mathcal S,
\]

то продолжение на алгебру задается формулой

\[
\mu(A)=\sum_{k=1}^n\mu(S_k).
\]

Главный вопрос на этом шаге - корректность: сумма не должна зависеть от
выбранного разбиения \(A\).

Продолжение на \(\sigma\)-алгебру строится через внешнюю меру:

\[
\mu^*(E)=
\inf\left\{
\sum_{n=1}^{\infty}\mu(A_n):
E\subset\bigcup_{n=1}^{\infty}A_n,\ A_n\in\mathcal A
\right\}.
\]

Из внешней меры выделяют измеримые множества по критерию Каратеодори:

\[
\mu^*(B)=\mu^*(B\cap E)+\mu^*(B\setminus E)
\]

для всех \(B\subset\Omega\). Эти множества образуют \(\sigma\)-алгебру,
а \(\mu^*\) на ней становится счетно-аддитивной мерой. При
\(\sigma\)-конечности продолжение на \(\sigma(\mathcal A)\) единственно.

Мера Лебега получается как продолжение объема полуоткрытых
прямоугольников:

\[
\lambda_d\left(\prod_{j=1}^d[a_j,b_j)\right)
=
\prod_{j=1}^d(b_j-a_j).
\]

На прямой это просто длина:

\[
\lambda([a,b))=b-a.
\]

\subsection{2. Полный
ответ}\label{ux43fux43eux43bux43dux44bux439-ux43eux442ux432ux435ux442}

\subsubsection{3.1. Введение и необходимость продолжения
меры}\label{ux432ux432ux435ux434ux435ux43dux438ux435-ux438-ux43dux435ux43eux431ux445ux43eux434ux438ux43cux43eux441ux442ux44c-ux43fux440ux43eux434ux43eux43bux436ux435ux43dux438ux44f-ux43cux435ux440ux44b}

В билете 1 были введены системы множеств: полуалгебры, алгебры и
\(\sigma\)-алгебры. В билете 2 вероятность была определена как
счетно-аддитивная мера на \(\sigma\)-алгебре. Теперь нужно объяснить,
откуда такие меры берутся.

Обычно меру невозможно удобно задать сразу на всей \(\sigma\)-алгебре.
Например, для меры Лебега естественно начать не со всех борелевских
множеств, а с полуинтервалов на прямой:

\[
\lambda([a,b))=b-a.
\]

В \(\mathbb R^d\) естественно начать с полуоткрытых прямоугольников:

\[
I=\prod_{j=1}^d[a_j,b_j),
\qquad
\lambda_d(I)=\prod_{j=1}^d(b_j-a_j).
\]

Эти множества простые: их легко пересекать, разрезать и считать их
длину, площадь или объем. Но для вероятности, интеграла Лебега и
случайных процессов нужны более сложные множества: счетные объединения,
счетные пересечения, пределы последовательностей множеств. Поэтому нужна
процедура продолжения меры.

\subsubsection{3.2. Продолжение с полуалгебры на
алгебру}\label{ux43fux440ux43eux434ux43eux43bux436ux435ux43dux438ux435-ux441-ux43fux43eux43bux443ux430ux43bux433ux435ux431ux440ux44b-ux43dux430-ux430ux43bux433ux435ux431ux440ux443}

Первый шаг - продолжение с полуалгебры \(\mathcal S\) на алгебру
\(\mathcal A(\mathcal S)\). Алгебра \(\mathcal A(\mathcal S)\) состоит
из конечных дизъюнктных объединений элементов \(\mathcal S\):

\[
\mathcal A(\mathcal S)=
\left\{
\bigsqcup_{k=1}^n S_k:\ S_k\in\mathcal S,\ n<\infty
\right\}.
\]

Если \(A\in\mathcal A(\mathcal S)\), то \(A\) представимо в виде

\[
A=\bigsqcup_{k=1}^n S_k.
\]

Естественное определение:

\[
\mu(A)=\sum_{k=1}^n\mu(S_k).
\]

Нужно доказать, что если у \(A\) есть другое разбиение, то сумма
получится та же. Это делается через общее измельчение двух разбиений.
Если

\[
A=\bigsqcup_{i=1}^m S_i=\bigsqcup_{j=1}^n T_j,
\]

то элементы \(S_i\cap T_j\) дают общее разбиение. Так как полуалгебра
замкнута относительно пересечений, эти множества снова лежат в
\(\mathcal S\). Аддитивность на \(\mathcal S\) дает равенство сумм.

\subsubsection{\texorpdfstring{3.3. Продолжение на \(\sigma\)-алгебру:
внешняя
мера}{3.3. Продолжение на \textbackslash sigma-алгебру: внешняя мера}}\label{ux43fux440ux43eux434ux43eux43bux436ux435ux43dux438ux435-ux43dux430-sigma-ux430ux43bux433ux435ux431ux440ux443-ux432ux43dux435ux448ux43dux44fux44f-ux43cux435ux440ux430}

Второй шаг - продолжение с алгебры на \(\sigma\)-алгебру. Пусть \(\mu\)
- предмера на алгебре \(\mathcal A\), то есть счетно-аддитивная мера на
\(\mathcal A\) в том смысле, что счетная аддитивность выполняется для
тех счетных дизъюнктных объединений, которые снова лежат в
\(\mathcal A\). Тогда строят внешнюю меру:

\[
\mu^*(E)=
\inf\left\{
\sum_{n=1}^{\infty}\mu(A_n):
E\subset\bigcup_{n=1}^{\infty}A_n,\ A_n\in\mathcal A
\right\}.
\]

Смысл: множество \(E\) покрывается счетным набором уже известных
измеримых множеств \(A_n\), а \(\mu^*(E)\) - наименьшая возможная цена
такого покрытия.

Внешняя мера определена на всех подмножествах \(\Omega\), но на всех
подмножествах она обычно не является мерой. Поэтому выделяют класс
хороших множеств.

\subsubsection{3.4. Условие измеримости по
Каратеодори}\label{ux443ux441ux43bux43eux432ux438ux435-ux438ux437ux43cux435ux440ux438ux43cux43eux441ux442ux438-ux43fux43e-ux43aux430ux440ux430ux442ux435ux43eux434ux43eux440ux438}

Множество \(E\) называется измеримым по Каратеодори, если для любого
\(B\subset\Omega\):

\[
\mu^*(B)=\mu^*(B\cap E)+\mu^*(B\setminus E).
\]

Это означает, что \(E\) корректно разрезает любое множество \(B\) на две
части без потери и без лишней массы. Теорема Каратеодори утверждает, что
все такие \(E\) образуют \(\sigma\)-алгебру, а \(\mu^*\) на ней является
счетно-аддитивной мерой. Кроме того, эта \(\sigma\)-алгебра содержит
исходную алгебру \(\mathcal A\), и новая мера совпадает с исходной
\(\mu\) на \(\mathcal A\).

\subsubsection{\texorpdfstring{3.5. Единственность продолжения и
\(\sigma\)-конечность}{3.5. Единственность продолжения и \textbackslash sigma-конечность}}\label{ux435ux434ux438ux43dux441ux442ux432ux435ux43dux43dux43eux441ux442ux44c-ux43fux440ux43eux434ux43eux43bux436ux435ux43dux438ux44f-ux438-sigma-ux43aux43eux43dux435ux447ux43dux43eux441ux442ux44c}

Если исходная мера \(\sigma\)-конечна, то продолжение на
\(\sigma(\mathcal A)\) единственно. \(\sigma\)-конечность означает, что
все пространство можно покрыть счетным числом множеств конечной меры:

\[
\Omega=\bigcup_{n=1}^{\infty}E_n,
\qquad
\mu(E_n)<\infty.
\]

Для меры Лебега это условие выполнено:

\[
\mathbb R^d=\bigcup_{n=1}^{\infty}[-n,n)^d,
\qquad
\lambda_d([-n,n)^d)=(2n)^d<\infty.
\]

\subsubsection{3.6. Резюме: построение меры
Лебега}\label{ux440ux435ux437ux44eux43cux435-ux43fux43eux441ux442ux440ux43eux435ux43dux438ux435-ux43cux435ux440ux44b-ux43bux435ux431ux435ux433ux430}

Мера Лебега строится именно по этой схеме. Сначала задают объем
полуоткрытого прямоугольника. Затем получают меру на алгебре
элементарных множеств, то есть конечных объединений прямоугольников.
После этого через внешнюю меру и критерий Каратеодори получают меру на
\(\sigma\)-алгебре. На борелевских множествах получается борелевская
мера Лебега, а после добавления всех подмножеств множеств меры ноль
получается полная мера Лебега.

\subsection{3. Основные
определения}\label{ux43eux441ux43dux43eux432ux43dux44bux435-ux43eux43fux440ux435ux434ux435ux43bux435ux43dux438ux44f}

\subsubsection{Продолжение
меры}\label{ux43fux440ux43eux434ux43eux43bux436ux435ux43dux438ux435-ux43cux435ux440ux44b}

Пусть \(\mathcal C\subset\mathcal D\). Мера \(\nu\) на \(\mathcal D\)
называется продолжением меры \(\mu\) с \(\mathcal C\), если

\[
\nu(A)=\mu(A),\qquad A\in\mathcal C.
\]

\subsubsection{Полуалгебра}\label{ux43fux43eux43bux443ux430ux43bux433ux435ux431ux440ux430}

Семейство \(\mathcal S\subset 2^\Omega\) такое, что
\(\varnothing\in\mathcal S\), пересечение двух множеств из
\(\mathcal S\) снова лежит в \(\mathcal S\), а разность двух множеств из
\(\mathcal S\) раскладывается в конечное дизъюнктное объединение
множеств из \(\mathcal S\):

\[
S\setminus T=\bigsqcup_{k=1}^n S_k,
\qquad S_k\in\mathcal S.
\]

\subsubsection{Алгебра, порожденная
полуалгеброй}\label{ux430ux43bux433ux435ux431ux440ux430-ux43fux43eux440ux43eux436ux434ux435ux43dux43dux430ux44f-ux43fux43eux43bux443ux430ux43bux433ux435ux431ux440ux43eux439}

\[
\mathcal A(\mathcal S)=
\left\{
\bigsqcup_{k=1}^n S_k:\ S_k\in\mathcal S,\ n<\infty
\right\}.
\]

Если \(\Omega\in\mathcal S\), то \(\mathcal A(\mathcal S)\) является
алгеброй.

\subsubsection{\texorpdfstring{\(\sigma\)-алгебра, порожденная системой
множеств}{\textbackslash sigma-алгебра, порожденная системой множеств}}\label{sigma-ux430ux43bux433ux435ux431ux440ux430-ux43fux43eux440ux43eux436ux434ux435ux43dux43dux430ux44f-ux441ux438ux441ux442ux435ux43cux43eux439-ux43cux43dux43eux436ux435ux441ux442ux432}

\[
\sigma(\mathcal C)=
\bigcap\{\mathcal F:\ \mathcal F\text{ - }\sigma\text{-алгебра},\ \mathcal C\subset\mathcal F\}.
\]

Это наименьшая \(\sigma\)-алгебра, содержащая \(\mathcal C\).

\subsubsection{Предмера на
алгебре}\label{ux43fux440ux435ux434ux43cux435ux440ux430-ux43dux430-ux430ux43bux433ux435ux431ux440ux435}

Функция \(\mu:\mathcal A\to[0,\infty]\) называется предмерой, если
\(\mu(\varnothing)=0\) и для любых попарно непересекающихся
\(A_1,A_2,\ldots\in\mathcal A\), таких что
\(\bigsqcup_{n=1}^{\infty}A_n\in\mathcal A\), выполнено

\[
\mu\left(\bigsqcup_{n=1}^{\infty}A_n\right)
=
\sum_{n=1}^{\infty}\mu(A_n).
\]

\subsubsection{\texorpdfstring{\(\sigma\)-конечность}{\textbackslash sigma-конечность}}\label{sigma-ux43aux43eux43dux435ux447ux43dux43eux441ux442ux44c}

Мера \(\mu\) называется \(\sigma\)-конечной, если существуют множества
\(E_n\) конечной меры такие, что

\[
\Omega=\bigcup_{n=1}^{\infty}E_n,
\qquad
\mu(E_n)<\infty.
\]

\subsubsection{Внешняя
мера}\label{ux432ux43dux435ux448ux43dux44fux44f-ux43cux435ux440ux430}

Функция \(\mu^*:2^\Omega\to[0,\infty]\), для которой:

\begin{enumerate}
\def\labelenumi{\arabic{enumi}.}
\tightlist
\item
  \(\mu^*(\varnothing)=0\);
\item
  \(A\subset B\Rightarrow\mu^*(A)\le\mu^*(B)\);
\item
  для любых \(A_1,A_2,\ldots\subset\Omega\):
\end{enumerate}

\[
\mu^*\left(\bigcup_{n=1}^{\infty}A_n\right)
\le
\sum_{n=1}^{\infty}\mu^*(A_n).
\]

\subsubsection{Измеримость по
Каратеодори}\label{ux438ux437ux43cux435ux440ux438ux43cux43eux441ux442ux44c-ux43fux43e-ux43aux430ux440ux430ux442ux435ux43eux434ux43eux440ux438}

Множество \(E\subset\Omega\) называется \(\mu^*\)-измеримым, если для
любого \(B\subset\Omega\):

\[
\mu^*(B)=\mu^*(B\cap E)+\mu^*(B\setminus E).
\]

\subsubsection{Мера
Лебега}\label{ux43cux435ux440ux430-ux43bux435ux431ux435ux433ux430}

Полная мера \(\lambda_d\) на \(\mathbb R^d\), полученная продолжением
объема полуоткрытых прямоугольников:

\[
\lambda_d\left(\prod_{j=1}^d[a_j,b_j)\right)
=
\prod_{j=1}^d(b_j-a_j).
\]

\subsection{4. Основные теоремы и
свойства}\label{ux43eux441ux43dux43eux432ux43dux44bux435-ux442ux435ux43eux440ux435ux43cux44b-ux438-ux441ux432ux43eux439ux441ux442ux432ux430}

\subsubsection{Теорема 1. Продолжение с полуалгебры на
алгебру}\label{ux442ux435ux43eux440ux435ux43cux430-1.-ux43fux440ux43eux434ux43eux43bux436ux435ux43dux438ux435-ux441-ux43fux43eux43bux443ux430ux43bux433ux435ux431ux440ux44b-ux43dux430-ux430ux43bux433ux435ux431ux440ux443}

Пусть \(\mathcal S\) - полуалгебра, \(\Omega\in\mathcal S\), и на
\(\mathcal S\) задана неотрицательная функция \(\mu\) с
\(\mu(\varnothing)=0\), согласованная с конечными дизъюнктными
разложениями: если

\[
S=\bigsqcup_{k=1}^n S_k,\qquad S,S_k\in\mathcal S,
\]

то

\[
\mu(S)=\sum_{k=1}^n\mu(S_k).
\]

Тогда на \(\mathcal A(\mathcal S)\) можно определить

\[
\mu\left(\bigsqcup_{k=1}^n S_k\right)
=
\sum_{k=1}^n\mu(S_k).
\]

Это определение корректно и задает аддитивное продолжение на алгебру.

\subsubsection{Доказательство}\label{ux434ux43eux43aux430ux437ux430ux442ux435ux43bux44cux441ux442ux432ux43e}

Пусть

\[
A=\bigsqcup_{i=1}^m S_i
=
\bigsqcup_{j=1}^n T_j,
\qquad S_i,T_j\in\mathcal S.
\]

Тогда для каждого \(i\):

\[
S_i=S_i\cap A
=
\bigsqcup_{j=1}^n(S_i\cap T_j).
\]

Так как \(\mathcal S\) замкнута относительно пересечений, все
\(S_i\cap T_j\) лежат в \(\mathcal S\). По аддитивности:

\[
\mu(S_i)=\sum_{j=1}^n\mu(S_i\cap T_j).
\]

Суммируем по \(i\):

\[
\sum_{i=1}^m\mu(S_i)
=
\sum_{i=1}^m\sum_{j=1}^n\mu(S_i\cap T_j).
\]

С другой стороны, для каждого \(j\):

\[
T_j=\bigsqcup_{i=1}^m(S_i\cap T_j),
\]

поэтому

\[
\sum_{j=1}^n\mu(T_j)
=
\sum_{j=1}^n\sum_{i=1}^m\mu(S_i\cap T_j).
\]

Двойные суммы одинаковы, значит значения по двум разбиениям совпадают.
\(\square\)

\subsubsection{Теорема 2. Теорема продолжения
меры}\label{ux442ux435ux43eux440ux435ux43cux430-2.-ux442ux435ux43eux440ux435ux43cux430-ux43fux440ux43eux434ux43eux43bux436ux435ux43dux438ux44f-ux43cux435ux440ux44b}

Пусть \(\mu\) - предмера на алгебре \(\mathcal A\). Тогда существует
мера \(\nu\) на \(\sigma\)-алгебре \(\mathcal M\), содержащей
\(\sigma(\mathcal A)\), такая что

\[
\nu(A)=\mu(A),\qquad A\in\mathcal A.
\]

Если \(\mu\) \(\sigma\)-конечна, то продолжение на
\(\sigma(\mathcal A)\) единственно.

\subsubsection{Схема
доказательства}\label{ux441ux445ux435ux43cux430-ux434ux43eux43aux430ux437ux430ux442ux435ux43bux44cux441ux442ux432ux430}

Строится внешняя мера

\[
\mu^*(E)=
\inf\left\{
\sum_{n=1}^{\infty}\mu(A_n):
E\subset\bigcup_{n=1}^{\infty}A_n,\ A_n\in\mathcal A
\right\}.
\]

Затем берутся все множества \(E\), удовлетворяющие критерию Каратеодори:

\[
\mu^*(B)=\mu^*(B\cap E)+\mu^*(B\setminus E),
\qquad B\subset\Omega.
\]

Доказывается, что они образуют \(\sigma\)-алгебру, что \(\mu^*\) на ней
счетно-аддитивна, что все множества из \(\mathcal A\) измеримы и что
\(\mu^*=\mu\) на \(\mathcal A\). Единственность при
\(\sigma\)-конечности доказывается разбиением пространства на счетное
объединение множеств конечной меры и применением теоремы единственности
на каждом таком куске.

\subsubsection{Теорема 3. Построение меры
Лебега}\label{ux442ux435ux43eux440ux435ux43cux430-3.-ux43fux43eux441ux442ux440ux43eux435ux43dux438ux435-ux43cux435ux440ux44b-ux43bux435ux431ux435ux433ux430}

На полуалгебре

\[
\mathcal S_d=
\left\{
\prod_{j=1}^d[a_j,b_j):\ a_j\le b_j
\right\}
\]

формула

\[
\lambda_d\left(\prod_{j=1}^d[a_j,b_j)\right)
=
\prod_{j=1}^d(b_j-a_j)
\]

задает объем. Он продолжается на алгебру элементарных множеств, затем на
\(\mathcal B(\mathbb R^d)\), а после пополнения дает полную меру Лебега.

\subsubsection{Обоснование}\label{ux43eux431ux43eux441ux43dux43eux432ux430ux43dux438ux435}

Полуоткрытые прямоугольники образуют полуалгебру. Объем аддитивен при
конечном разрезании прямоугольников. Значит его можно продолжить на
алгебру элементарных множеств. Эта мера \(\sigma\)-конечна, потому что

\[
\mathbb R^d=\bigcup_{n=1}^{\infty}[-n,n)^d,
\qquad
\lambda_d([-n,n)^d)<\infty.
\]

По теореме продолжения получается единственная борелевская мера,
совпадающая с объемом на прямоугольниках. Ее пополнение и называется
мерой Лебега.

\subsubsection{Свойство 4. Внешняя мера
Лебега}\label{ux441ux432ux43eux439ux441ux442ux432ux43e-4.-ux432ux43dux435ux448ux43dux44fux44f-ux43cux435ux440ux430-ux43bux435ux431ux435ux433ux430}

Для \(E\subset\mathbb R^d\):

\[
\lambda_d^*(E)
=
\inf\left\{
\sum_{n=1}^{\infty}\lambda_d(I_n):
E\subset\bigcup_{n=1}^{\infty}I_n,\ I_n\in\mathcal S_d
\right\}.
\]

Если \(E\) лебегово измеримо, то

\[
\lambda_d(E)=\lambda_d^*(E).
\]

\subsubsection{Свойство 5. Инвариантность меры Лебега относительно
сдвига}\label{ux441ux432ux43eux439ux441ux442ux432ux43e-5.-ux438ux43dux432ux430ux440ux438ux430ux43dux442ux43dux43eux441ux442ux44c-ux43cux435ux440ux44b-ux43bux435ux431ux435ux433ux430-ux43eux442ux43dux43eux441ux438ux442ux435ux43bux44cux43dux43e-ux441ux434ux432ux438ux433ux430}

Если \(E\) лебегово измеримо и \(h\in\mathbb R^d\), то

\[
\lambda_d(E+h)=\lambda_d(E),
\qquad
E+h=\{x+h:\ x\in E\}.
\]

Идея доказательства: сдвиг прямоугольника не меняет его объем, поэтому
не меняет суммарные объемы покрытий. Значит внешняя мера сохраняется, а
критерий Каратеодори переносит измеримость на сдвинутое множество.

\subsubsection{Свойство 6. Точки и счетные множества имеют меру
ноль}\label{ux441ux432ux43eux439ux441ux442ux432ux43e-6.-ux442ux43eux447ux43aux438-ux438-ux441ux447ux435ux442ux43dux44bux435-ux43cux43dux43eux436ux435ux441ux442ux432ux430-ux438ux43cux435ux44eux442-ux43cux435ux440ux443-ux43dux43eux43bux44c}

Для любой точки \(x\in\mathbb R^d\):

\[
\lambda_d(\{x\})=0.
\]

\subsubsection{Доказательство}\label{ux434ux43eux43aux430ux437ux430ux442ux435ux43bux44cux441ux442ux432ux43e-1}

Точку можно покрыть кубом со стороной \(\varepsilon>0\), поэтому

\[
\lambda_d^*(\{x\})\le\varepsilon^d.
\]

Так как \(\varepsilon\) можно брать сколь угодно малым, получаем
\(\lambda_d(\{x\})=0\).

Если \(E=\{x_1,x_2,\ldots\}\) счетно, то по счетной субаддитивности:

\[
\lambda_d(E)
\le
\sum_{n=1}^{\infty}\lambda_d(\{x_n\})
=0.
\]

Значит \(\lambda_d(E)=0\). \(\square\)

\subsection{5. Примеры}\label{ux43fux440ux438ux43cux435ux440ux44b}

\subsubsection{Пример 1. Длина конечного объединения
полуинтервалов}\label{ux43fux440ux438ux43cux435ux440-1.-ux434ux43bux438ux43dux430-ux43aux43eux43dux435ux447ux43dux43eux433ux43e-ux43eux431ux44aux435ux434ux438ux43dux435ux43dux438ux44f-ux43fux43eux43bux443ux438ux43dux442ux435ux440ux432ux430ux43bux43eux432}

Пусть

\[
A=[0,1)\sqcup[2,5).
\]

Тогда

\[
\lambda(A)=1+3=4.
\]

Если разбить иначе:

\[
A=[0,\tfrac12)\sqcup[\tfrac12,1)\sqcup[2,3)\sqcup[3,5),
\]

то

\[
\lambda(A)=\frac12+\frac12+1+2=4.
\]

Это иллюстрирует независимость от разбиения.

\subsubsection{Пример 2. Открытые и закрытые
интервалы}\label{ux43fux440ux438ux43cux435ux440-2.-ux43eux442ux43aux440ux44bux442ux44bux435-ux438-ux437ux430ux43aux440ux44bux442ux44bux435-ux438ux43dux442ux435ux440ux432ux430ux43bux44b}

После продолжения меры:

\[
\lambda((a,b))=\lambda([a,b))=\lambda((a,b])=\lambda([a,b])=b-a.
\]

Концы не влияют на меру, потому что точки имеют меру \(0\).

\subsubsection{Пример 3. Рациональные числа на
отрезке}\label{ux43fux440ux438ux43cux435ux440-3.-ux440ux430ux446ux438ux43eux43dux430ux43bux44cux43dux44bux435-ux447ux438ux441ux43bux430-ux43dux430-ux43eux442ux440ux435ux437ux43aux435}

Множество \(\mathbb Q\cap[0,1]\) счетно, поэтому

\[
\lambda(\mathbb Q\cap[0,1])=0.
\]

Но оно всюду плотно в \([0,1]\). Значит мера Лебега измеряет длину, а не
топологическую плотность.

\subsubsection{Пример 4. Иррациональные числа на
отрезке}\label{ux43fux440ux438ux43cux435ux440-4.-ux438ux440ux440ux430ux446ux438ux43eux43dux430ux43bux44cux43dux44bux435-ux447ux438ux441ux43bux430-ux43dux430-ux43eux442ux440ux435ux437ux43aux435}

Так как

\[
[0,1]=(\mathbb Q\cap[0,1])\sqcup([0,1]\setminus\mathbb Q),
\]

то

\[
\lambda([0,1]\setminus\mathbb Q)=1.
\]

\subsubsection{Пример 5. Канторово
множество}\label{ux43fux440ux438ux43cux435ux440-5.-ux43aux430ux43dux442ux43eux440ux43eux432ux43e-ux43cux43dux43eux436ux435ux441ux442ux432ux43e}

Классическое канторово множество несчетно, но имеет меру \(0\).
Суммарная длина удаленных интервалов равна

\[
\sum_{n=1}^{\infty}\frac{2^{n-1}}{3^n}
=
\frac13\sum_{n=0}^{\infty}\left(\frac23\right)^n
=1.
\]

Из отрезка длины \(1\) удалена суммарная длина \(1\), поэтому оставшееся
канторово множество имеет меру \(0\).

\subsection{6. Доказательства и
обоснования}\label{ux434ux43eux43aux430ux437ux430ux442ux435ux43bux44cux441ux442ux432ux430-ux438-ux43eux431ux43eux441ux43dux43eux432ux430ux43dux438ux44f}

\textbf{Почему продолжение на алгебру задается суммой.}

Если \(A=\bigsqcup_{k=1}^n S_k\), то любая аддитивная мера обязана
удовлетворять

\[
\mu(A)=\sum_{k=1}^n\mu(S_k).
\]

Иначе она нарушала бы аддитивность. Поэтому формула единственно
возможная. Сложность только в том, чтобы доказать независимость от
разбиения.

\textbf{Почему нужна внешняя мера.}

Алгебра закрыта только относительно конечных операций. Но
\(\sigma\)-алгебра содержит результаты счетных операций. Внешняя мера
назначает размер любому множеству через счетные покрытия уже известными
множествами:

\[
E\subset\bigcup_{n=1}^{\infty}A_n.
\]

Инфимум сумм \(\sum_n\mu(A_n)\) дает наилучшую оценку размера \(E\)
сверху.

\textbf{Почему не все множества сразу измеримы.}

Внешняя мера субаддитивна на всех подмножествах, но аддитивность может
ломаться. Поэтому оставляют только те множества \(E\), которые правильно
разрезают любое \(B\):

\[
\mu^*(B)=\mu^*(B\cap E)+\mu^*(B\setminus E).
\]

Это и есть критерий Каратеодори.

\textbf{Почему \(\sigma\)-конечность важна для единственности.}

Если пространство покрыто счетным числом множеств конечной меры, то
совпадение двух продолжений можно проверять на каждом конечном куске, а
затем склеивать по счетному покрытию. Без \(\sigma\)-конечности
продолжение может быть не единственным.

\textbf{Почему мера Лебега продолжает геометрический объем.}

На прямоугольниках объем известен. При конечном разрезании
прямоугольника сумма объемов частей равна объему целого. Это дает меру
на полуалгебре и алгебре элементарных множеств. Теорема продолжения
переносит эту меру на борелевские множества, а пополнение дает полную
меру Лебега.

\subsection{7. Что написать на
доске}\label{ux447ux442ux43e-ux43dux430ux43fux438ux441ux430ux442ux44c-ux43dux430-ux434ux43eux441ux43aux435}

\begin{enumerate}
\def\labelenumi{\arabic{enumi}.}
\tightlist
\item
  Схема:
\end{enumerate}

\[
\mathcal S\to\mathcal A(\mathcal S)\to\sigma(\mathcal S).
\]

\begin{enumerate}
\def\labelenumi{\arabic{enumi}.}
\setcounter{enumi}{1}
\tightlist
\item
  Алгебра из полуалгебры:
\end{enumerate}

\[
\mathcal A(\mathcal S)=
\left\{
\bigsqcup_{k=1}^n S_k:\ S_k\in\mathcal S
\right\}.
\]

\begin{enumerate}
\def\labelenumi{\arabic{enumi}.}
\setcounter{enumi}{2}
\tightlist
\item
  Продолжение на алгебру:
\end{enumerate}

\[
\mu\left(\bigsqcup_{k=1}^n S_k\right)
=
\sum_{k=1}^n\mu(S_k).
\]

\begin{enumerate}
\def\labelenumi{\arabic{enumi}.}
\setcounter{enumi}{3}
\tightlist
\item
  Внешняя мера:
\end{enumerate}

\[
\mu^*(E)=
\inf\left\{
\sum_{n=1}^{\infty}\mu(A_n):
E\subset\bigcup_{n=1}^{\infty}A_n,\ A_n\in\mathcal A
\right\}.
\]

\begin{enumerate}
\def\labelenumi{\arabic{enumi}.}
\setcounter{enumi}{4}
\tightlist
\item
  Критерий Каратеодори:
\end{enumerate}

\[
\mu^*(B)=\mu^*(B\cap E)+\mu^*(B\setminus E).
\]

\begin{enumerate}
\def\labelenumi{\arabic{enumi}.}
\setcounter{enumi}{5}
\tightlist
\item
  Мера Лебега:
\end{enumerate}

\[
\lambda_d\left(\prod_{j=1}^d[a_j,b_j)\right)
=
\prod_{j=1}^d(b_j-a_j).
\]

\begin{enumerate}
\def\labelenumi{\arabic{enumi}.}
\setcounter{enumi}{6}
\tightlist
\item
  \(\sigma\)-конечность:
\end{enumerate}

\[
\mathbb R^d=\bigcup_{n=1}^{\infty}[-n,n)^d.
\]

\subsection{8. Типичные дополнительные
вопросы}\label{ux442ux438ux43fux438ux447ux43dux44bux435-ux434ux43eux43fux43eux43bux43dux438ux442ux435ux43bux44cux43dux44bux435-ux432ux43eux43fux440ux43eux441ux44b}

\textbf{Вопрос. Почему нельзя сразу пользоваться формулой конечной суммы
на \(\sigma(\mathcal S)\)?}

Потому что элементы \(\sigma(\mathcal S)\) могут не иметь конечного
разбиения на элементы \(\mathcal S\). Они возникают через счетные
объединения, пересечения и пределы.

\textbf{Вопрос. Что такое предмера?}

Это счетно-аддитивная мера на алгебре, еще не продолженная на
\(\sigma\)-алгебру.

\textbf{Вопрос. Что дает \(\sigma\)-конечность?}

Единственность продолжения на порожденную \(\sigma\)-алгебру.

\textbf{Вопрос. Внешняя мера определена на всех подмножествах. Значит,
все подмножества измеримы?}

Нет. Внешняя мера есть на всех подмножествах, но счетно-аддитивной мерой
она становится только на \(\sigma\)-алгебре измеримых по Каратеодори
множеств.

\textbf{Вопрос. Чем борелевская мера Лебега отличается от полной меры
Лебега?}

Борелевская мера задана на \(\mathcal B(\mathbb R^d)\). Полная мера
Лебега дополнительно содержит все подмножества борелевских множеств меры
ноль.

\textbf{Вопрос. Почему \(\lambda([a,b])=b-a\), хотя исходно задавали
\([a,b)\)?}

Потому что

\[
[a,b]=[a,b)\cup\{b\},
\]

а точка имеет меру ноль.

\textbf{Вопрос. Может ли несчетное множество иметь меру ноль?}

Да. Пример - канторово множество.

\subsection{9. Типичные
ошибки}\label{ux442ux438ux43fux438ux447ux43dux44bux435-ux43eux448ux438ux431ux43aux438}

\begin{enumerate}
\def\labelenumi{\arabic{enumi}.}
\tightlist
\item
  Путать классы множеств:
\end{enumerate}

\[
\mathcal S\subset\mathcal A(\mathcal S)\subset\sigma(\mathcal S).
\]

\begin{enumerate}
\def\labelenumi{\arabic{enumi}.}
\setcounter{enumi}{1}
\tightlist
\item
  Забывать, что формула
\end{enumerate}

\[
\mu\left(\bigsqcup_{k=1}^nS_k\right)=\sum_{k=1}^n\mu(S_k)
\]

требует попарной непересекаемости.

\begin{enumerate}
\def\labelenumi{\arabic{enumi}.}
\setcounter{enumi}{2}
\item
  Не проверять независимость от разбиения.
\item
  Считать внешнюю меру полноценной мерой на всех подмножествах.
\item
  Забывать \(\sigma\)-конечность в утверждении об единственности
  продолжения.
\item
  Путать борелевскую меру Лебега и полную меру Лебега.
\item
  Думать, что множество меры ноль обязательно счетно.
\item
  Использовать старое обозначение дополнения через верхний индекс. В
  этих конспектах дополнение пишется чертой сверху:
\end{enumerate}

\[
\overline{A}=\Omega\setminus A.
\]

\subsection{10. Связи с другими
билетами}\label{ux441ux432ux44fux437ux438-ux441-ux434ux440ux443ux433ux438ux43cux438-ux431ux438ux43bux435ux442ux430ux43cux438}

\textbf{Билет 01.} Полуалгебры, алгебры и \(\sigma\)-алгебры нужны для
построения мер через продолжение.

\textbf{Билет 02.} Вероятностная мера - частный случай меры. Мера по
функции распределения строится той же идеей продолжения.

\textbf{Билет 04.} Борелевские множества и пополнение меры продолжают
тему меры Лебега.

\textbf{Билет 05.} Измеримые функции определяются относительно
\(\sigma\)-алгебр, полученных такими построениями.

\textbf{Билет 07.} Интеграл Лебега строится по мере.

\textbf{Билет 11.} Замена переменной в интеграле использует свойства
меры Лебега.

\textbf{Билет 15 и Билет 16.} Построение мер на пространствах
последовательностей является более сложной версией идеи продолжения.

\subsection{11. Минимум на
удовлетворительно}\label{ux43cux438ux43dux438ux43cux443ux43c-ux43dux430-ux443ux434ux43eux432ux43bux435ux442ux432ux43eux440ux438ux442ux435ux43bux44cux43dux43e}

Нужно уметь сказать:

\begin{enumerate}
\def\labelenumi{\arabic{enumi}.}
\tightlist
\item
  Меру задают на простом классе множеств и продолжают:
\end{enumerate}

\[
\mathcal S\to\mathcal A(\mathcal S)\to\sigma(\mathcal S).
\]

\begin{enumerate}
\def\labelenumi{\arabic{enumi}.}
\setcounter{enumi}{1}
\tightlist
\item
  На алгебре конечных дизъюнктных объединений:
\end{enumerate}

\[
\mu\left(\bigsqcup_{k=1}^nS_k\right)=\sum_{k=1}^n\mu(S_k).
\]

\begin{enumerate}
\def\labelenumi{\arabic{enumi}.}
\setcounter{enumi}{2}
\item
  На \(\sigma\)-алгебру продолжение строится через внешнюю меру и
  критерий Каратеодори.
\item
  Мера Лебега продолжает длину:
\end{enumerate}

\[
\lambda([a,b))=b-a.
\]

\begin{enumerate}
\def\labelenumi{\arabic{enumi}.}
\setcounter{enumi}{4}
\tightlist
\item
  В \(\mathbb R^d\) она продолжает объем прямоугольника:
\end{enumerate}

\[
\lambda_d\left(\prod_{j=1}^d[a_j,b_j)\right)
=
\prod_{j=1}^d(b_j-a_j).
\]

\subsection{12. Минимум на
хорошо}\label{ux43cux438ux43dux438ux43cux443ux43c-ux43dux430-ux445ux43eux440ux43eux448ux43e}

Кроме минимума на удовлетворительно, нужно:

\begin{enumerate}
\def\labelenumi{\arabic{enumi}.}
\tightlist
\item
  Объяснить корректность продолжения на алгебру через общее измельчение
  разбиений.
\item
  Сформулировать внешнюю меру:
\end{enumerate}

\[
\mu^*(E)=
\inf\left\{
\sum_{n=1}^{\infty}\mu(A_n):
E\subset\bigcup_{n=1}^{\infty}A_n
\right\}.
\]

\begin{enumerate}
\def\labelenumi{\arabic{enumi}.}
\setcounter{enumi}{2}
\tightlist
\item
  Сформулировать критерий Каратеодори:
\end{enumerate}

\[
\mu^*(B)=\mu^*(B\cap E)+\mu^*(B\setminus E).
\]

\begin{enumerate}
\def\labelenumi{\arabic{enumi}.}
\setcounter{enumi}{3}
\tightlist
\item
  Сказать, что при \(\sigma\)-конечности продолжение единственно.
\item
  Привести примеры: длина интервала, мера точки, мера
  \(\mathbb Q\cap[0,1]\).
\end{enumerate}

\subsection{13. Что добавить для
отлично}\label{ux447ux442ux43e-ux434ux43eux431ux430ux432ux438ux442ux44c-ux434ux43bux44f-ux43eux442ux43bux438ux447ux43dux43e}

Для сильного ответа стоит добавить:

\begin{enumerate}
\def\labelenumi{\arabic{enumi}.}
\tightlist
\item
  Различие между мерой на полуалгебре, предмерой на алгебре и мерой на
  \(\sigma\)-алгебре.
\item
  Подробную схему доказательства теоремы продолжения через внешнюю меру.
\item
  Объяснение, почему внешняя мера не является мерой на всех
  подмножествах.
\item
  Различие между борелевской мерой Лебега и полной мерой Лебега.
\item
  Доказательство, что точка и любое счетное множество имеют меру ноль.
\item
  Пример канторова множества как несчетного множества меры ноль.
\item
  Упоминание, что без \(\sigma\)-конечности единственность продолжения
  может нарушаться.
\end{enumerate}

\subsection{14. Устная версия ответа на 5
минут}\label{ux443ux441ux442ux43dux430ux44f-ux432ux435ux440ux441ux438ux44f-ux43eux442ux432ux435ux442ux430-ux43dux430-5-ux43cux438ux43dux443ux442}

Тема билета - как построить меру на сложной \(\sigma\)-алгебре, если
сначала она задана только на простых множествах. Схема:

\[
\mathcal S\to\mathcal A(\mathcal S)\to\sigma(\mathcal S).
\]

Сначала есть полуалгебра \(\mathcal S\), например полуоткрытые интервалы
или прямоугольники. Из нее строится алгебра \(\mathcal A(\mathcal S)\)
конечных дизъюнктных объединений. Если

\[
A=\bigsqcup_{k=1}^nS_k,
\]

то кладем

\[
\mu(A)=\sum_{k=1}^n\mu(S_k).
\]

Главное доказать, что это не зависит от разбиения. Если есть два
разбиения \(A=\bigsqcup_iS_i=\bigsqcup_jT_j\), берем пересечения
\(S_i\cap T_j\). Они дают общее измельчение, и по аддитивности суммы
совпадают.

Дальше строится внешняя мера:

\[
\mu^*(E)=
\inf\left\{
\sum_{n=1}^{\infty}\mu(A_n):
E\subset\bigcup_{n=1}^{\infty}A_n,\ A_n\in\mathcal A
\right\}.
\]

Она задает размер любого множества через наилучшие счетные покрытия.
Затем по критерию Каратеодори выбираются измеримые множества:

\[
\mu^*(B)=\mu^*(B\cap E)+\mu^*(B\setminus E).
\]

Они образуют \(\sigma\)-алгебру, и на ней \(\mu^*\) становится мерой.
При \(\sigma\)-конечности продолжение на порожденную \(\sigma\)-алгебру
единственно.

Мера Лебега - главный пример. На полуоткрытых прямоугольниках задаем

\[
\lambda_d\left(\prod_{j=1}^d[a_j,b_j)\right)
=
\prod_{j=1}^d(b_j-a_j).
\]

Затем продолжаем на алгебру элементарных множеств, потом на борелевскую
\(\sigma\)-алгебру и после пополнения получаем полную меру Лебега. На
прямой:

\[
\lambda([a,b))=b-a.
\]

Точки и счетные множества имеют меру ноль, например
\(\lambda(\mathbb Q\cap[0,1])=0\), а \(\lambda([0,1])=1\).

\subsection{15. Если осталось 2 минуты перед
экзаменом}\label{ux435ux441ux43bux438-ux43eux441ux442ux430ux43bux43eux441ux44c-2-ux43cux438ux43dux443ux442ux44b-ux43fux435ux440ux435ux434-ux44dux43aux437ux430ux43cux435ux43dux43eux43c}

Запомнить:

\[
\mathcal S\to\mathcal A(\mathcal S)\to\sigma(\mathcal S).
\]

На алгебре:

\[
\mu\left(\bigsqcup_{k=1}^nS_k\right)=\sum_{k=1}^n\mu(S_k).
\]

Внешняя мера:

\[
\mu^*(E)=
\inf\left\{
\sum_n\mu(A_n):
E\subset\bigcup_n A_n
\right\}.
\]

Каратеодори:

\[
\mu^*(B)=\mu^*(B\cap E)+\mu^*(B\setminus E).
\]

Мера Лебега:

\[
\lambda([a,b))=b-a,
\qquad
\lambda_d\left(\prod_{j=1}^d[a_j,b_j)\right)
=
\prod_{j=1}^d(b_j-a_j).
\]

Главные ловушки: забыть независимость от разбиения, спутать внешнюю меру
с мерой на всех подмножествах и забыть \(\sigma\)-конечность для
единственности.

\newpage

\end{document}
